\section*{Extended abstract}
\textbf{Study question:} How do reference composition, infertility subtype and outcome/denominator definitions affect the interpretation of PCOS and endometriosis infertility comparisons?

\textbf{Summary answer:} Reference-country shares explained a substantial part of the modeled China--G20 PCOS contrast change, while a separate survey showed that diagnosis-associated estimates varied with definitions and denominators.

\textbf{What is known already:} GBD studies describe infertility burden, trends and subtypes. Age standardization does not fix country composition within a pooled reference, and survey definitions are not interchangeable.

\textbf{Study design, size, duration:} Ecological secondary analysis of GBD 2023, 1990--2023, covering 19 sovereign G20 countries and seven female age groups, plus a post hoc cross-sectional NSFG 2022--2023 analysis with 4,855 respondents aged 20--49 at interview.

\textbf{Participants/materials, setting, methods:} China was compared with 18 other countries using pooling within age, restricted-age standardization, subtype accounting and symmetric endpoint decomposition, with annual-chained sensitivity. NSFG analysis used final weights, 20 strata and 80 clusters, preserving joint Taylor covariance. Independent R and Python implementations checked estimates and covariance.

\textbf{Main results and the role of chance:} The PCOS-related infertility contrast narrowed by 151.54 per 100,000; reference composition contributed 88.65 (58.5\%), compared with 87.11 (57.5\%) under annual chaining. Secondary infertility accounted for 69.4\%. These GBD quantities lack propagated uncertainty. Survey prevalence ratios for broad fertility problems versus 12-month infertility in unions were 2.65 (95\% design CI 2.30--3.04) versus 1.77 (1.22--2.57) for PCOS, and 1.78 (1.40--2.25) versus 0.84 (0.50--1.43) for endometriosis. The paired within-union PCOS definition shift in prevalence differences was 17.15 percentage points (10.24--24.06).

\textbf{Limitations, reasons for caution:} No jointly aligned GBD draws were obtained. Survey associations are unadjusted and partly sparse; the narrow endometriosis interval crosses one. Definitions and populations differ, and NSFG is not a proven held-out GBD validation sample. Several analyses were post hoc; endometriosis terminal-year instability is unresolved.

\textbf{Wider implications of the findings:} Comparisons should distinguish subtype, reference composition, definitions and denominators. Neither convergence nor definition-sensitive associations demonstrate improved care, disease causation or a protective endometriosis effect.

\textbf{Study funding/disclosures:} Actual human-author funding and disclosures remain to be verified; none are assumed.

\textbf{Trial registration number:} Not applicable; no interventional trial or external preregistration is claimed.

\textbf{Key words:} infertility; PCOS; endometriosis; reference population; survey design
